\chapter{Egzotyczne struktury gładkie na \texorpdfstring{$\R^4$}{R4}}
\label{ch:egzotyczne-r4}
\index{struktura gładka!egzotyczna na $\R^4$}
\index{przestrzeń euklidesowa!egzotyczna $\R^4$}

Z poprzednich rozdziałów znamy już różnicę między
rozmaitością topologiczną a gładką. Klasyfikacja
Freedmana z Rozdziału~\ref{ch:freedman-klasyfikacja}
rozstrzyga pewne pytania o homeomorfizmy, lecz nie
wyznacza atlasu gładkiego. W tym rozdziale wykażemy,
że na przestrzeni topologicznej $\R^4$ istnieje
nieskończenie wiele struktur gładkich, które nie są
ze sobą dyfeomorficzne.

Dowód jest konstrukcją indukcyjną Gompfa. Jego
wejściami są wyniki Freedmana o uchwytach Cassona,
omówionych w Rozdziale~\ref{ch:uchwyty-cassona},
przeszkoda gładka Donaldsona oraz dwa głębokie
lematy Gompfa. Sformułujemy te wejścia jawnie.
Następnie samodzielnie przeprowadzimy krok
indukcyjny i sprawdzimy, dlaczego zwarte podzbiory
rozróżniają otrzymane typy gładkie.

\section{Jak rozróżniać struktury gładkie}
\label{sec:egzotyczne-r4-pojecia}

\begin{definicja}[Egzotyczna $\R^4$]
\label{def:egzotyczna-r4}
Gładka, spójna rozmaitość $R$ jest \emph{egzotyczną
$\R^4$}, jeżeli jest homeomorficzna z topologiczną
$\R^4$, ale nie jest dyfeomorficzna ze standardową
gładką $\R^4$. Dwie struktury gładkie na tej samej
topologicznej $\R^4$ uznajemy za równoważne, gdy
istnieje między nimi dyfeomorfizm.
\end{definicja}

Intuicja: homeomorfizm pozwala używać tych samych
ciągłych współrzędnych globalnych, lecz nie musi
przenosić różniczkowalności. Przykładem odniesienia
jest standardowy atlas na $\R^4$; sama zamiana
globalnych współrzędnych przez dyfeomorfizm
nie daje nowego typu gładkiego. Egzotyczny przykład
otrzymamy dopiero po skonstruowaniu niezanurzalnego
zwartego podzbioru.

W całym dowodzie przestrzenie są zorientowane.
Przez \emph{zanurzenie} rozumiemy gładkie zanurzenie
zachowujące orientację, chyba że wyraźnie napiszemy
inaczej. Odwzorowania topologiczne będziemy tak
nazywać. Rozróżnienie to jest istotne: rezultat
Gompfa najpierw odróżnia typy \emph{zorientowane}.

Niech $K$ będzie zwartą gładką podrozmaitością
czterowymiarową z brzegiem, zanurzoną w $R$.
Jeżeli $K$ nie da się zanurzyć gładko w $R'$, to
nie istnieje gładkie zanurzenie $R\hookrightarrow R'$:
jego ograniczenie do $K$ dałoby zabronione
zanurzenie. Tę prostą obserwację zastosujemy do
ciągu starannie dobranych $K$.

\section{Ślady splotów i sumy przy końcu}
\label{sec:egzotyczne-r4-sploty}

Przypomnijmy terminologię uchwytową z
Definicji~\ref{def:freedman-framed-link}.

\begin{definicja}[Topologicznie plastrowalny splot]
\label{def:egzotyczne-r4-slice}
Splot $L=L_1\sqcup\cdots\sqcup L_s\subset S^3=\partial B^4$
nazywamy \emph{topologicznie plastrowalnym}, jeżeli
jego składowe ograniczają parami rozłączne,
lokalnie płasko zanurzone dyski $D_i\subset B^4$.
Przez $M_L$ oznaczamy zwartą gładką
czterowymiarową rozmaitość otrzymaną z $B^4$
przez dołączenie do $L$ uchwytów $2$ z obramowaniem
zerowym, mierzonym względem obramowania Seiferta.
\end{definicja}

Lokalna płaskość oznacza, że para $(B^4,D_i)$
w otoczeniu punktu dysku wygląda topologicznie
jak $(\R^4,\R^2)$ lub odpowiedni model przy brzegu.
To warunek potrzebny do istnienia topologicznego
otoczenia dysku o postaci $D^2\times D^2$.
Intuicyjnie, rdzenie dołączonych uchwytów w $M_L$
zastępują dyski plastra; ich obramowania muszą
się zgadzać. Przykład stanowi splot złożony z $s$
niezaczepionych okręgów: ograniczają zwykłe,
rozłączne gładkie dyski w $B^4$, a $M_L$ jest
brzegową sumą $s$ kopii $S^2\times D^2$.
Warunek topologicznej plastrowalności nie orzeka,
czy istnieją odpowiednie \emph{gładkie} dyski.

\begin{definicja}[Ogolona $\R^4$ i suma przy końcu]
\label{def:egzotyczne-r4-shaved}
Gładką rozmaitość $R$ homeomorficzną z $\R^4$
nazywamy \emph{ogoloną} (ang. \emph{shaved}),
jeżeli istnieje jej
gładkie zanurzenie w pewnej gładkiej rozmaitości
$M$ na wnętrze lokalnie płasko zanurzonej
topologicznej zamkniętej kuli $B\subset M$, przy
czym niepusty otwarty fragment $\partial B$ jest
gładką podrozmaitością $M$.

Jeżeli $R=\operatorname{int}B$ oraz
$R'=\operatorname{int}B'$ są tak przedstawione,
umieszczamy $B$ i $B'$ w jednej rozmaitości,
wykonując sumę spójną rozmaitości otaczających
w kulach rozłącznych z $B$ i $B'$.
Wybieramy łuk łączący gładkie fragmenty ich
brzegów. Po dołączeniu cienkiego otoczenia łuku
otrzymujemy brzegową sumę kul; wnętrze nazywamy
\emph{sumą przy końcu} i zapisujemy
$R\mathbin{\natural_\infty}R'$.
\end{definicja}

Cały brzeg $B$ nie musi być gładki: właśnie jego
brak gładkości może przechowywać egzotyczną
strukturę wnętrza. Gładki fragment wystarcza,
by wykonać brzegową sumę. Przykładem ogolonej
rozmaitości jest standardowe $\R^4$, czyli
wnętrze zwykłej gładkiej kuli $B^4\subset S^4$.
Brzegowa suma dwóch zwykłych kul jest znowu
zwykłą kulą, więc odpowiadająca jej suma przy
końcu jest standardową $\R^4$.

\begin{stwierdzenie}[Własności używanej sumy przy końcu]
\label{prop:egzotyczne-r4-suma}
Niech $R,R'$ będą ogolonymi gładkimi rozmaitościami
homeomorficznymi z $\R^4$. Sumę przy końcu można
wybrać tak, by była ogolona i homeomorficzna z
$\R^4$, a zadane zwarte podzbiory $K\subset R$ oraz
$K'\subset R'$ zanurzały się w niej gładko i
rozłącznie. Ponadto $R$ i $R'$ zanurzają się w
tej sumie z zachowaniem orientacji.
\end{stwierdzenie}
\begin{proof}
Po wspomnianej sumie spójnej rozmaitości otaczających
wybierzmy gładkie fragmenty brzegów kul $B,B'$
oraz łuk pomiędzy nimi, nieprzecinający zadanych
zwartych podzbiorów. Otoczenie łuku jest
jednym brzegowym uchwytem $1$ łączącym dwie
topologiczne kule. Ich brzegowa suma jest ponownie
topologiczną kulą $4$-wymiarową: topologicznie
można wyprostować łuk i oglądać dwie kule
połączone cienką rurą. Brzeg pozostaje lokalnie
płaski, a na fragmentach odległych od końców
łuku pozostaje gładki. Wnętrze tej kuli jest
zatem homeomorficzne z $\R^4$ i spełnia definicję
ogolenia. Wnętrza $B$ i $B'$ wchodzą w nie jako
rozłączne gładkie otwarte podrozmaitości, co daje
żądane zanurzenia wraz z $K$ i $K'$.
\end{proof}

Używamy tu \emph{konkretnej} sumy wybranych
modeli. Pełna niezależność jej typu gładkiego
od wyboru promieni i łuku nie jest potrzebna
w poniższej indukcji; jej dowód znajduje się
w dodatku do pracy Gompfa podanej niżej.

\section{Gładka przeszkoda i trzy wejścia zewnętrzne}
\label{sec:egzotyczne-r4-wejscia}

Mechanizm niezanurzalności bierze początek z
Twierdzenia~\ref{thm:donaldson-diagonalizacja}.
Jego założenia, źródło i rachunek energii równań ASD
omówiliśmy w Rozdziale~\ref{ch:donaldson}.
Samo twierdzenie o
diagonalizacji nie podaje jeszcze zwartego zbioru,
który będzie potrzebny w indukcji.

\begin{twierdzenie}[Gompf: stała przeszkoda zwarta;
wynik zewnętrzny]
\label{thm:egzotyczne-r4-X}
Istnieje zwarta, spójna, zorientowana gładka
czterowymiarowa rozmaitość $X$ o spójnym brzegu
oraz ogolona rozmaitość $R_1$ homeomorficzna
z $\R^4$ takie, że $X$ zanurza się gładko w $R_1$,
lecz $X$ nie zanurza się gładko z zachowaniem
orientacji w żadnej zamkniętej, gładkiej,
jednospójnej czterowymiarowej rozmaitości
o ujemnie określonej formie przecięcia.
\end{twierdzenie}

Taką parę $(X,R_1)$ wybiera
\href{https://doi.org/10.4310/jdg/1214439566}
{R.~E.~Gompf, \emph{An infinite set of exotic
$\R^4$'s}, J.~Differential Geom. 21 (1985),
\S\,1, s.~285--286}, korzystając ze swojej wcześniejszej
\href{https://doi.org/10.4310/jdg/1214437666}
{pracy \emph{Three exotic $\R^4$'s and other
anomalies} (1983)}. Konstrukcja początkowego
$R_1$ wykorzystuje topologiczne uchwyty Cassona
i przeszkodę Donaldsona. Gompf podaje także
konkretną możliwość wyboru $X$: ślad $M_J$
zerowo obramowanego węzła $J$ otrzymanego przez
cztery podwojenia Whiteheada z ujemnymi klamrami
lewoskrętnego węzła trójlistnego. Uzasadnienie
niezanurzalności tego szczególnego $M_J$ wymaga
dodatkowych wyników o uchwytach Cassona; nie
wyprowadzamy go tu z samego diagramu węzła.

\begin{stwierdzenie}[Proste konsekwencje przeszkody]
\label{prop:egzotyczne-r4-X-konsekwencje}
Rozmaitość $X$ z
Twierdzenia~\ref{thm:egzotyczne-r4-X} nie zanurza
się gładko z zachowaniem orientacji w żadnej
gładkiej homotopijnej $4$-sferze. W szczególności
nie zanurza się w standardowej $\R^4$.
\end{stwierdzenie}
\begin{proof}
Gdyby $X\hookrightarrow\Sigma$, gdzie $\Sigma$
jest gładką homotopijną $4$-sferą, wykonajmy
sumę spójną $\Sigma\#\overline{\C P^2}$
w kuli rozłącznej z obrazem $X$. Rozmaitość
ta jest zamknięta, gładka i jednospójna.
Ponieważ $H_2(\Sigma;\Z)=0$, jej forma przecięcia
jest równa $\langle-1\rangle$. Zanurzenie $X$
pozostałoby w tej rozmaitości, co przeczy
Twierdzeniu~\ref{thm:egzotyczne-r4-X}.
Standardowa $\R^4$ zanurza się gładko w $S^4$,
więc także do niej $X$ nie może się zanurzyć.
\end{proof}

\begin{twierdzenie}[Gompf: realizacja śladu splotu;
wynik zewnętrzny]
\label{thm:egzotyczne-r4-slad}
Dla każdego topologicznie plastrowalnego splotu
$L\subset S^3$ istnieje ogolona gładka
rozmaitość $R_L$ homeomorficzna z $\R^4$,
w której $M_L$ zanurza się gładko z zachowaniem
orientacji.
\end{twierdzenie}

Jest to \href{https://doi.org/10.4310/jdg/1214439566}
{lemat~1.1 Gompfa (1985)}. Topologiczne dyski
dają najpierw topologiczne zanurzenie śladu
uchwytowego w $\R^4$. Gładkość w pobliżu
$M_L$ nadaje się przez przeniesienie jego atlasu;
przedłużenie wygładzenia na dopełnienie
oraz wybór gładkiego fragmentu brzegu wymagają
twierdzeń o wygładzaniu końców. W dowodzie
Gompf odwołuje się w tym miejscu do
\href{https://doi.org/10.4310/jdg/1214437139}
{F.~Quinna, \emph{Ends of maps, III: dimensions
4 and 5}, J.~Differential Geom. 17 (1982)}.

\begin{twierdzenie}[Gompf: nowa przeszkoda dla
każdego etapu; wynik zewnętrzny]
\label{thm:egzotyczne-r4-lemat}
Niech $R$ będzie dowolną ogoloną, zorientowaną
gładką rozmaitością homeomorficzną z $\R^4$.
Istnieje topologicznie plastrowalny splot
$L\subset S^3$ taki, że rozłączna suma
$X\sqcup M_L$ nie ma gładkiego zanurzenia
zachowującego orientację w $R$. Tutaj $X$
jest \emph{tym samym} zwartym blokiem co w
Twierdzeniu~\ref{thm:egzotyczne-r4-X}.
\end{twierdzenie}

To dokładnie \href{https://doi.org/10.4310/jdg/1214439566}
{lemat~1.2 Gompfa (1985)}. Hipoteza ogolenia
jest w nim istotna: pozwala umieścić $R$
jako wnętrze topologicznej kuli w gładkiej
rozmaitości otaczającej. Nie zastępujemy jej
samym założeniem $R\cong_{\mathrm{TOP}}\R^4$.

\begin{uwaga}[Geometria lematu~1.2]
\label{rem:egzotyczne-r4-mechanizm}
Warto zobaczyć, co wymusza zakaz zanurzenia.
Gompf umieszcza ogoloną kulę $B$ w zamkniętej,
jednospójnej, gładkiej rozmaitości $M$ o
hiperbolicznej formie przecięcia. Poza $B$
reprezentuje hiperboliczną bazę przy użyciu
uchwytów Cassona. Pięć pierwszych pięter
każdego uchwytu jest zwartą wieżą. Rdzenie
szóstego piętra są jeszcze dyskami immersyjnymi;
Gompf przenosi ich samoprzecięcia ku kuli bazowej
i z pozostałych gładkich dysków buduje kule $B_i$.
Przecięcia $\partial B_i$ z topologicznymi
sferami reprezentującymi bazę tworzą sploty $L_i$;
ich rozłączna suma w rozłącznych kulach
trójwymiarowych daje poszukiwany $L$.
Splot jest topologicznie plastrowalny, ponieważ
odpowiednie uchwyty Cassona są topologicznie
otwartymi uchwytami $2$; jest to dokładnie
Twierdzenie~\ref{thm:casson-homeomorphism}.

Gdyby $X\sqcup M_L$ zanurzyło się w $R\subset B$,
gładkie rdzenie uchwytów $2$ w $M_L$ dałyby
dyski zastępujące topologiczne fragmenty sfer
zbudowanych poza $B$. Powstałyby pary gładkich
sfer $S_i,T_i$, z
$S_i^2=T_i^2=0$ i $S_i\cdot T_i=1$,
rozłączne między różnymi parami, reprezentujące
całe $H_2(M;\Z)$. Każda para jest konfiguracją
$S^2\vee S^2$: sfery przecinają się w jednym
punkcie. Jej regularne otoczenie jest
$(S^2\times S^2)\setminus\operatorname{int}B^4$.
Wycięcie tych otoczeń i wklejenie zwykłych
kul $4$ daje gładką homotopijną $4$-sferę,
w której $X$ nadal się zanurza. Jest to sprzeczne
ze Stwierdzeniem~\ref{prop:egzotyczne-r4-X-konsekwencje}.
Istotnie, otoczenia mają brzeg $S^3$ i
reprezentują całe $H_2(M;\Z)$; po ich wycięciu
oraz zaklejeniu brzegu kulami $B^4$ twierdzenia
van Kampena i Mayera--Vietorisa dają odpowiednio
$\pi_1=0$ i $H_2=0$. Dualność Poincarégo
wyznacza pozostałe grupy homologii zamkniętej
rozmaitości, więc ma ona typ homotopijny $S^4$.

Ten opis wskazuje miejsce każdej przeszkody,
ale nie jest dowodem lematu~1.2: konstrukcja
zanurzonych uchwytów Cassona, kontrola obramowań
i rozłączność dysków zajmują istotną część
oryginalnego argumentu Gompfa. Jego głębokie
wejścia pochodzą z
\href{https://doi.org/10.4310/jdg/1214437136}
{M.~H.~Freedmana, \emph{The topology of
four-dimensional manifolds} (1982)} oraz
\href{https://doi.org/10.4310/jdg/1214437139}
{Quinna (1982)}.
\end{uwaga}

\section{Indukcja Gompfa i zwarte świadki}
\label{sec:egzotyczne-r4-indukcja}

\begin{twierdzenie}[Nieskończenie wiele struktur
gładkich na $\R^4$]
\label{thm:egzotyczne-r4-nieskonczenie}
Istnieje nieskończenie wiele parami niedyfeomorficznych
gładkich rozmaitości homeomorficznych z $\R^4$.
Można je wybrać spośród rozmaitości Gompfa
$R_0,R_1,R_2,\ldots$, przy czym $R_0$ jest
standardową $\R^4$, a każda $R_n$ dla $n\geq1$
jest egzotyczna.
\end{twierdzenie}
\begin{proof}[Dowód z twierdzeń zewnętrznych
\ref{thm:egzotyczne-r4-X},
\ref{thm:egzotyczne-r4-slad} i
\ref{thm:egzotyczne-r4-lemat}]
Połóżmy $R_0=\R^4$ ze standardową orientacją.
Rozmaitość $R_1$ i jej zwarty podzbiór $X$
pochodzą z Twierdzenia~\ref{thm:egzotyczne-r4-X}.
Stwierdzenie~\ref{prop:egzotyczne-r4-X-konsekwencje}
pokazuje, że nie ma zanurzenia $R_1\hookrightarrow R_0$.

Załóżmy, że ogolona $R_{n-1}$, zawierająca
gładko $X$, została już zbudowana, gdzie $n\geq2$.
Z Twierdzenia~\ref{thm:egzotyczne-r4-lemat}
wybieramy splot $L_n$ taki, że zwarta
rozłączna suma
\begin{equation}\label{eq:egzotyczne-r4-Kn}
 K_n:=X\sqcup M_{L_n}
 \quad\text{nie zanurza się gładko w }R_{n-1}.
\end{equation}
Twierdzenie~\ref{thm:egzotyczne-r4-slad}
daje ogoloną $R_{L_n}$ z zanurzonym
$M_{L_n}$. Tworzymy
\begin{equation}\label{eq:egzotyczne-r4-Rn}
 R_n:=R_{n-1}\mathbin{\natural_\infty}R_{L_n}.
\end{equation}
Łuk sumowania wybieramy z dala od $X$ i
$M_{L_n}$. Stwierdzenie~\ref{prop:egzotyczne-r4-suma}
zapewnia, że $R_n$ jest ogolona,
homeomorficzna z $\R^4$, zawiera $R_{n-1}$
oraz zawiera rozłącznie zanurzone oba składniki
$K_n$. W szczególności w $R_n$ nadal jest $X$,
więc krok można powtarzać.

Gdyby istniało gładkie zanurzenie zachowujące
orientację $R_n\hookrightarrow R_{n-1}$,
jego ograniczenie do $K_n$ przeczyłoby
\eqref{eq:egzotyczne-r4-Kn}. Zatem
$R_n\not\hookrightarrow R_{n-1}$.
Dla $m<n$ mamy natomiast zanurzenia
$R_m\hookrightarrow R_{m+1}\hookrightarrow
\cdots\hookrightarrow R_{n-1}$. Gdyby
$R_n\hookrightarrow R_m$, ich złożenie
dałoby wykluczone zanurzenie
$R_n\hookrightarrow R_{n-1}$. Otrzymaliśmy
więc nieskończenie wiele parami różnych
\emph{zorientowanych} typów dyfeomorfizmu.
Każda $R_n$ dla $n\geq1$ jest egzotyczna:
gdyby była dyfeomorficzna ze standardową
$R_0$ z zachowaniem orientacji, mielibyśmy
zanurzenie $R_n\hookrightarrow R_0$. Jeśli
dyfeomorfizm odwracałby orientację, można
złożyć go z odbiciem standardowej $\R^4$;
wynik byłby dyfeomorfizmem zachowującym
orientację, co prowadzi do tej samej sprzeczności.

Trzeba jeszcze sprawdzić tezę bez warunku
zachowania orientacji. Jeden niezorientowany
typ dyfeomorfizmu ma najwyżej dwa typy
zorientowane: odpowiadają one dwóm wyborom
orientacji. Gdyby wśród $R_n$ występowało
tylko skończenie wiele typów niezorientowanych,
byłoby także tylko skończenie wiele typów
zorientowanych. Jest ich nieskończenie wiele,
więc istnieje nieskończony podciąg $R_{n_j}$
parami niedyfeomorficzny nawet po dopuszczeniu
odwrócenia orientacji.
Możemy włączyć do tego podciągu $R_0$.

Na koniec wybierzmy dla każdego $j$ homeomorfizm
$h_j:\R^4\to R_{n_j}$ i przenieśmy atlas gładki
$R_{n_j}$ na tę samą topologiczną przestrzeń
$\R^4$. Dyfeomorfizm między dwoma przeniesionymi
atlasami dałby dyfeomorfizm między odpowiednimi
$R_{n_j}$. Takiego dyfeomorfizmu nie ma, więc
otrzymane struktury gładkie są parami nierównoważne.
\end{proof}

W dowodzie widać dokładnie rolę zwartego świadka:
$K_n$ \emph{zanurza się} w $R_n$, ale \emph{nie zanurza
się} w $R_{n-1}$. Sama odmienność splotów $L_n$
albo diagramów uchwytowych nie wystarczyłaby do
rozróżnienia rozmaitości. Silniejsze twierdzenie
Gompfa, \href{https://doi.org/10.4310/jdg/1214439566}
{twierdzenie~1.3}, opisuje pełny porządek zanurzeń
zorientowanych: $R_m\hookrightarrow R_n$ dokładnie
wtedy, gdy $m\leq n$.
